\documentclass{article}
\input{boilerplate}

\begin{document}


\newcommand{\lectureTitle}{Reinforcement Learning for Large-Scale Auto-Bidding}
\newcommand{\lectureHeader}{RL for Large-Scale Auto-Bidding}
\newcommand{\lectureDate}{Due: May 8th, 10:30 PM}

\textsc{COS435 / ECE433: Introduction to RL} 
\vspace{1em} \hfill \lectureDate

\maketitle

\paragraph{Authors.}($\alpha$-$\beta$ order) Ruirong Feng (\texttt{rrfeng@princeton.edu}), Ophelia Fulton (\texttt{of4642@princeton.edu}),
Xinya (Claire) Gu (\texttt{xg3412@princeton.edu}),
Bonnie Liu (\texttt{bl0919@princeton.edu}),
Yiqi (Andrew) Liu (\texttt{andrew.liu@princeton.edu})

\paragraph{Link to GitHub Repo.} \url{https://github.com/liuyiqiandrew/AuctionNet}
% \section*{Instructions}
%
% This is a rough template guide for your final project report. If you do not have these exact sections in your final report, or do not follow this formatting, do not worry -- prioritize the contents over the exact section breakdown here.
%
% A few things to keep in mind:
% \begin{itemize}
%     \item {\bf The main body of your final report must be $\leq 8$ pages.} If your current report is longer than $8$ pages, please push finer details to an Appendix or cut down on the current text. You should {\it probably} aim to have at least 4 pages of content, but again, we won't automatically take off points if the length is too short. 
%     \item {\bf There is no minimum length for the final report.} Please make sure to cover the rubric items in the final project guidelines.
%     \item {\bf The minimum allowed font size if 10 pt. as it is in this template for readability.} There is no restriction on the specific font you choose. Because there is no minimum length, you can choose to make your font-size 12 pt., but please keep everything reasonable-looking (i.e. no 16 pt. situation). 
%     \item {\bf Signpost.} Please do some signposting, even if you don't use the section breakdown here -- this is to ensure that we don't miss anything within your report. It is okay to even, i.e., include headings for your paragraphs:
%     \\\\{\bf What is the problem? } Insert text here. \\
%
%     but we encourage you to write this report like a realistic NeurIPS, ICLR, ICML, etc. paper where content is wrapped into the main sections. Again, we will evaluate the reports on {\bf content} over the exact formatting and structure. 
% \end{itemize}

\section{Introduction}

% Introduce the problem. Please cover the following points:
% \begin{itemize}
%     \item Motivation: What is the problem?
%     \item Motivation: Why is it interesting and important?
%     \item Motivation/relationship with prior work: Why is it hard? (E.g., why do naive approaches fail?)
%     \item Relationship with prior work: Why hasn't it been solved before? (Or, what's wrong with previous proposed solutions? How does mine differ?)
%     \item Methods and results: What are the key components of my approach and results? Also include any specific limitations.
% \end{itemize}

Auto-bidding is the control layer between an advertiser's campaign objective and
the real-time auction market. An advertiser aims to maximize conversions, under
a limited budget and the realized cost per acquisition (CPA).
AuctionNet turns this problem into a public benchmark built from large-scale
auction logs: at each decision point, the bidder chooses a scalar multiplier
that determines bids for a batch of impressions, competes against other
advertisers, and receives reward only after auction outcomes and conversions are
realized \citep{su2024auctionnet,xu2024autobidding}. The task is therefore a
natural sequential decision problem, but not a clean textbook one. The agent
observes only a summary of the market, the relevant constraints operate over an
episode, and the best action depends on how much budget and CPA slack should be
saved for future traffic.

Static or myopic bidding rules miss this tradeoff. A policy that bids
aggressively on every high-probability impression can exhaust its budget before
later opportunities arrive; a policy that only protects CPA can leave valuable
conversions unclaimed. Standard on-policy actor-critic training is a natural
baseline because it learns directly from simulator feedback, but the naive
version leaves several failure modes exposed: the CPA penalty is delayed, the
observation hides market-regime changes, and uniform rollouts spend the same
gradient budget on easy and hard periods.

Recent alternatives make larger changes to the solution concept.
Dynamic-knapsack methods decompose long-horizon advertising into budget
allocation and policy optimization, but depend on a knapsack structure and
greedy selection rule that is not the same as learning a scalar bid multiplier
inside AuctionNet's replay simulator \citep{hao2020dynamic}. DiffBid avoids the
Markov assumption by training a conditional diffusion model over whole
trajectories, but this shifts the problem toward offline trajectory generation
from historical bidding logs rather than online policy improvement
\citep{guo2024generative}. Oracle imitation learning is especially strong on
AuctionNet, but its supervision comes from a hindsight oracle that sees the full
campaign traffic and solves a nonlinear multiple-choice knapsack problem before
training the student policy \citep{chiappa2025autobidding}.

These methods identify what vanilla online RL does not specify on its own: how
delayed CPA feedback should shape pacing, which market statistics should be
exposed to a partially observed policy, how recent history should enter the
state, and how training updates should be allocated across heterogeneous
periods. We study these factors while keeping the AuctionNet scalar action,
replay simulator, and evaluation protocol fixed. Starting from the provided
actor-critic policy, we formulate the task as a POMDP and use the same baseline
pipeline as the reference point for all extensions
\citep{cai2017realtime,wu2018budget,he2021unified}.

Our experiments vary the components of the online RL pipeline rather than
replacing the benchmark with a different supervision problem. At the observation
level, we compare compact benchmark features, a richer state used by stronger
competition agents, and additional market-regime signals. At the training level,
reward shaping and score-weighted period sampling test whether the policy
benefits from denser constraint feedback and more informative rollouts. At the
policy level, a recurrent encoder tests whether learned history summaries
improve over hand-crafted rolling statistics. We also evaluate large language
models as bidding policies by fine-tuning Qwen3-8B with GRPO and LoRA adapters.
Across these experiments, improvements in state design account for the largest
gains, while the pretrained LLM policy shows faster early adaptation but a lower
ceiling than the strongest non-LLM actor-critic policy in our runs. The resulting
study provides a compact training recipe for constrained decision making under
partial information and an empirical probe of how generic pretrained policies
behave in numerical constrained optimization with delayed rewards.

% \section{Related Work}
% Literature review here; a rough sketch: 
% \begin{itemize}
%   \item Works that put forward the auto-bidding problem, to show why it is important
%   \item Different ramifications of solutions to auto-bidding
%   \item Review of RL methods and why they are natural and compelling for this problem
% \end{itemize}

\section{Methods}
% we may also need one big ML conference style diagram here 
%
% NOTE: put hyperparameter tuning process in the Appendix, talk high-level things here.

\subsection{Problem Setup and PPO Baseline}
We study the AuctionNet auto-bidding environment from the NeurIPS 2024 competition on auto-bidding under uncertainty \citep{xu2024autobidding}.
The environment follows the real-time auction setting used in recent
auto-bidding work \citep{chiappa2025autobidding}. Each
episode corresponds to one advertiser in one delivery period. The advertiser is
given a total budget \(B\), a target cost per acquisition
\(\mathrm{CPA}_{\mathrm{target}}\), and a horizon of \(T=48\) bidding ticks.
At tick \(t\), the simulator presents a batch \(\mathcal{I}_t\) of impression
opportunities, each with a predicted conversion probability
\(p_{t,i}\). The agent chooses a scalar log-bid multiplier
\(a_t \in [-10,10]\), which is applied to every impression in the batch:
\[
  b_{t,i} = \exp(a_t)\,p_{t,i}\,\mathrm{CPA}_{\mathrm{target}},
  \qquad i \in \mathcal{I}_t.
\]
The controlled bids are inserted into historical auctions against the other
47 advertisers, whose bids are fixed by the dataset. Slots and prices are then
determined by the auction simulator, spend is deducted from the remaining
budget, and conversions are sampled from the predicted values. The simulator
enforces the budget constraint, while the AuctionNet evaluation score applies a
quadratic penalty when the realized CPA exceeds the target. For a trajectory
\(\tau\), total conversions and cost are \(Y_T(\tau)\) and \(C_T(\tau)\),
respectively, and the realized CPA is
\(\mathrm{CPA}_T(\tau)=C_T(\tau)/\max\{Y_T(\tau),1\}\). The evaluation score is
\[
  R(\tau) =
  Y_T(\tau) \cdot
  \min\!\left\{1,
  \left(\frac{\mathrm{CPA}_{\mathrm{target}}}{\mathrm{CPA}_T(\tau)}\right)^2
  \right\}.
\]

We formulate the task as a partially observable Markov decision process
\((\mathcal{S}, \mathcal{A}, \Omega, \mathcal{P}, \mathcal{O},
\mathcal{R})\). The latent state \(s_t \in \mathcal{S}\) contains the
current campaign status, auction traffic, competitor bids, and future market
conditions that the agent cannot fully observe. The action space
\(\mathcal{A}\) is the scalar log-bid multiplier above. The observation
\(o_t \in \Omega\) is a feature vector summarizing the current tick, remaining
budget, target CPA, current impression statistics, and historical bidding
outcomes. After the agent samples \(a_t \sim \pi_\theta(\cdot \mid h_t)\), with
history \(h_t=(o_0,a_0,\ldots,o_t)\), the simulator transitions according to
\(\mathcal{P}(\cdot \mid s_t,a_t)\), emits a new observation through
\(\mathcal{O}\), and returns reward \(r_t=\mathcal{R}(s_t,a_t)\). In the
standard PPO baseline, this reward is the per-step CPA-clipped score computed
from the realized cost and conversions at tick \(t\); the terminal AuctionNet
score is used for evaluation. The expected return under policy \(\pi_\theta\)
is
\[
  J(\theta)=
  \mathbb{E}_{\pi_\theta}\!\left[
    \sum_{t=0}^{T-1} \gamma^t r_t
  \right],
\]
with discount factor \(\gamma\). Our baseline uses the 16-feature observation,
an MLP actor-critic policy with a diagonal Gaussian action distribution, and
running observation/reward normalization. It is trained on 20 periods of
trajectory data and evaluated on one held-out period.

For the baseline actor-critic policy, we optimize the standard clipped PPO
surrogate. Let \(\pi_{\theta_{\mathrm{old}}}\) be the rollout policy,
\(\rho_t(\theta)=\pi_\theta(a_t\mid o_t)/
\pi_{\theta_{\mathrm{old}}}(a_t\mid o_t)\), \(\hat A_t\) be the advantage
estimate, and \(\hat V_t\) be the bootstrapped return target. The PPO objective
is
\begin{equation}
\begin{split}
  \mathcal{L}_{\mathrm{PPO}}(\theta)
  =
  \mathbb{E}_t \Big[
    &\min\!\big(
      \rho_t(\theta)\hat A_t,\,
      \mathrm{clip}(\rho_t(\theta),1-\epsilon_c,1+\epsilon_c)\hat A_t
    \big) \\
    &\qquad\qquad\qquad- c_v\big(V_\theta(o_t)-\hat V_t\big)^2
    + c_H\,\mathcal{H}\!\left(\pi_\theta(\cdot\mid o_t)\right)
  \Big].
\end{split}
\end{equation}
Beyond the standard clipped PPO surrogate, the quadratic term fits the value
function used for advantage estimation, and the entropy term
\(\mathcal{H}(\pi_\theta)\) discourages premature collapse of the continuous
Gaussian policy.
This baseline fixes the auction simulator, scalar bid interface, and evaluation
protocol for all PPO-based extensions below.

\subsection{PPO Pipeline Extensions}
We treat the online PPO bidder as the reference pipeline and vary the components
most directly tied to constrained bidding under partial information: constraint
feedback, rollout allocation, temporal context, and market state. All variants share the AuctionNet scalar action interface and replay evaluation protocol, yielding a controlled ablation in which score differences are attributable to the modified component.

% @Andrew, Claire, Bonnie, Ophelia: write your extension here succintly --- motivation, main extension with formula (and explain your formula);
%
% Use the following environment to wrap up your extension:
%
% --- BEGIN ---
% \paragraph{Your extension.}
% Your writing
% --- END ---
\paragraph{Reward Shaping.}
The unconstrained PPO maximizes rewards but without mechanism to constrain CPA when the agent exceeds the target CPA. We added a fixed Lagrangian penalty coefficient $\lambda$ that penalizes CPA exceedance, applied per step rather than only at terminal:
\begin{align*}
  \mathcal{L}(\pi, \lambda) \;=\; \mathbb{E}\!\left[\sum_t \text{conv}_t\right] \;-\; \lambda \cdot \mathbb{E}\!\left[\bigl(\text{total\_cost} - \text{target\_CPA} \cdot \text{total\_conv}\bigr)_{+}\right]
\end{align*}
$\lambda$ here is a hyperparameter that we tune, and more specifically $\lambda \in \{.0001, .001, .01\}$. We calculated the per step overspend as:
\[
  \text{overspend}_t = \max\bigl(0,\; \text{cost}_t - \text{target\_cpa} \cdot \text{conv}_t\bigr)
\]
so the shaped reward $r^{\text{shaped}}_t = r_t - \lambda * \text{overspend}_t$, and when $\lambda = 0$, we have the original PPO. With such, we expect the policy learns to bid more carefully and follows CPA in training.

% --- BEGIN ---
\paragraph{Sample Reweighting.}
The default pipeline samples uniformly from training periods 7--26, spending equal
gradient budget on periods where the policy already performs well and those it struggles
with. We replace this with a score-weighted distribution that oversamples
high-performing (advertiser, period) pairs. For each pair~$j$ we compute the policy's
NeurIPS score and assign sampling weight
%
\begin{equation}
  w_j = \alpha \cdot \mathrm{softmax}(\mathrm{score}_j / T)
      + (1 - \alpha) \cdot \tfrac{1}{N},
\end{equation}
%
where $T$ controls concentration strength and $\alpha$ interpolates between reweighted
and uniform sampling. At each rollout, the worker draws a new pair according to~$w_j$.
We set $\alpha{=}0.9$ and $T{=}30$: the $(1{-}\alpha)$ uniform floor ensures every
period is visited occasionally, preventing the policy from forgetting rare regimes, while $T{=}30$ provides moderate concentration, which ensures the highest-scoring trajectories
receive several times the sampling probability of the lowest, but without collapsing
onto a small handful of easy campaigns. All other pipeline components remain unchanged.
% --- END ---

%%%%%%%%%%%%%%%%%%%%%%%%%%%%%%%%%%%%%%
\paragraph{Temporal Encoding.}

The original online PPO bidding policy observes each auction period as a
state vector $x_t \in \mathbb{R}^{d}$ and uses a feed-forward 
network to map the state vector to a bid coefficient.
The temporal extension alters the state representation for the policy
by using the $K$ most recent states as the input $X_t$ for each step.
If fewer than $K$ steps are available, the history is padded by repeating the earliest
available state.

The history $X_t$ is encoded into a hidden summary $h_t$ by a 
residue Gated Recurrent Unit (GRU) \cite{cho-etal-2014-learning}.
The hidden vector $h_t$ is concatenated with the latest state $x_t$ and 
passed through layer normalization to form the augmented PPO input $z_t$.
\begin{align}
    h_t = \mathrm{GRU}(X_t), \qquad
    z_t = \mathrm{LayerNorm}([h_t; x_t]).
\end{align}
This residue path preserves current-period bidding features, while
the GRU summarizes short-horizon trajectory information.

Unless otherwise specified, actor and critic use
separate temporal feature extractors, matching the intuition that action
selection and value estimation may attend to different temporal summaries.
A shared-extractor variant was also implemented to test whether the extra
capacity of separate GRUs is necessary or whether a single temporal
representation regularizes the actor-critic update.



\paragraph{Market-aware Feature Engineering.}
%%%%%%%%%%%%%%%%%%%%%%%%%%%%%%%%%%%%%%

\subsection{RL Training on LLM}
We test the performance of large language models (LLM) on the auto-bidding task
before and after RL training to explore whether a pretrained generic model can
improve final score, sample efficiency, and training stability of reinforcement learning. We use Qwen3-8B
as the base policy model. At each tick, the current auction state is serialized
into a short prompt containing budget, target CPA, market statistics, and historical
performance. The model emits one numeric action, treated as the same log-bid
multiplier $a_t$ used by the PPO baseline:
\[
  b_{t,i} = \exp(a_t)\, p_{t,i}\,\mathrm{CPA}_{\mathrm{target}}.
\]

We adopt Group Relative Policy Optimization (GRPO) as the training objective. For each prompt
\(x\), the rollout policy \(\pi_{\mathrm{old}}\) samples \(G\) trajectories
\(\{\tau_i\}_{i=1}^{G}\), receives terminal returns \(\{R_i\}_{i=1}^G\), and assigns each
trajectory the group-normalized advantage
\(\hat A_i=(R_i-\mu_R)/(\sigma_R+\varepsilon)\), where \(\mu_R\) and
\(\sigma_R\) are computed within the group. The optimization objective is the clipped surrogate
with a KL penalty to the reference policy:
\begin{equation}
\begin{split}
  \mathcal{L}_{\mathrm{GRPO}}(\theta)
  &=
  \mathbb{E}_{x,\{\tau_i\}}\Bigg[
  \frac{1}{G}\sum_{i=1}^{G}
  \min\!\Bigg(
    \frac{\pi_\theta(\tau_i\mid x)}
         {\pi_{\mathrm{old}}(\tau_i\mid x)}\hat A_i, \\
  &\qquad\qquad\qquad
    \mathrm{clip}\!\left(
      \frac{\pi_\theta(\tau_i\mid x)}
           {\pi_{\mathrm{old}}(\tau_i\mid x)},
      1-\epsilon_c,1+\epsilon_c
    \right)\hat A_i
  \Bigg)
  - \beta D_{\mathrm{KL}}(\pi_\theta\,\Vert\,\pi_{\mathrm{ref}})
  \Bigg].
\end{split}
\end{equation}
In the auto-bidding environment, a prompt \(x\) corresponds to one
\((\text{period},\text{advertiser})\) row. One trajectory \(\tau_i\) is a full
48-tick episode: at every tick, the environment state is rendered as the next
user message, the LLM emits the numeric action \(a_t\), and the simulator advances one
step. Chain-of-thought is disabled so that the assistant tokens are only the
action response. Each GRPO update samples four prompts and \(G=4\) trajectories
per prompt, giving 16 episodes whose terminal AuctionNet scores are the returns
\(R_i\) above. During training, we use LoRA adapters on Qwen3-8B's attention
and MLP projection modules to make policy updates feasible without
full-parameter fine-tuning. The experiment keeps
the AuctionNet dynamics and scalar bid interface fixed relative to the MLP
actor-critic baseline, isolating the effect of replacing that policy class with
an LLM.

\section{Results}

\begin{table}[h]
  \centering
  \small
  \caption{Held-out period-27 sweep score (mean $\pm$ s.e.m., $n{=}48$).
  Bold marks the best variant within each extension family.}
  \label{tab:ppo-extensions-summary}
  \begin{tabular}{@{}llccc@{}}
    \toprule
    Method & Variant & Obs & Steps & Sweep score ($n{=}48$) \\
    \midrule
    PPO baseline       & ---                                  & $16$ & $10$M & $20.14 \pm 2.25$ \\
    PPO baseline       & ---                                  & $60$ & $10$M & $23.23 \pm 2.02$ \\
    \midrule
    Reward shaping     & $\lambda = 10^{-2}$                  & $16$ & $10$M & $22.10 \pm 1.73$ \\
    Reward shaping     & $\lambda = 10^{-3}$                  & $16$ & $10$M & $25.33 \pm 2.18$ \\
    Reward shaping     & $\lambda = 10^{-4}$                  & $16$ & $10$M & $\mathbf{25.95 \pm 2.15}$ \\
    \midrule
    Sample reweighting & $\alpha{=}0.9,\,T{=}30$              & $16$ & $10$M & $24.64 \pm 2.23$ \\
    Sample reweighting & $\alpha{=}0.9,\,T{=}30$              & $60$ & $10$M & $\mathbf{28.59 \pm 2.12}$ \\
    \midrule
    Temporal encoding  & $K{=}4,\,H{=}64$, separate           & $16$ & $10$M & $24.78 \pm 2.29$ \\
    Temporal encoding  & $K{=}8,\,H{=}32$, shared             & $60$ & $10$M & $\mathbf{24.85 \pm 2.28}$ \\
    \midrule
    Market-aware       & $o_t^{(98)}$                         & $98$ & $5$M  & $\mathbf{24.87 \pm 2.21}$ \\
    Market-aware       & $o_t^{(98)}$                         & $98$ & $10$M & $24.85 \pm 2.09$ \\
    \bottomrule
  \end{tabular}
\end{table}

Table~\ref{tab:ppo-extensions-summary} compiles the held-out period-27 sweep
scores for the four PPO extensions on a common axis. The strongest
configuration per extension is sample reweighting on obs-60 ($28.59$),
followed by reward shaping at $\lambda{=}10^{-4}$ on obs-16 ($25.95$), the
$5$M-step market-aware run ($24.87$), and the best temporal encoder
($24.85$). All four extensions improve on the obs-16 baseline; sample
reweighting is the only one that also clears the stronger obs-60 baseline by
more than one s.e.m. The sections that follow report each extension in
detail.

\subsection{Temporal Encoding}
\begin{table}[h]
    \centering
    \small
    \caption{Held-out period-27 advertiser-sweep score for temporal encoding extension. 
    Entries report mean terminal scores.}
    \label{tab:temporal-eval}
    \begin{tabular}{@{}cccccc@{}}
        \toprule
        Features & $K$ & \multicolumn{2}{c}{$H=32$} & \multicolumn{2}{c}{$H=64$} \\
        \cmidrule(lr){3-4} \cmidrule(lr){5-6}
         &  & Separate & Shared & Separate & Shared \\
        \midrule
        16 & Baseline & \multicolumn{4}{c}{$20.14 \pm 2.25$} \\
        16 & 4 & $23.80 \pm 2.32$ & $23.58 \pm 2.29$ & $24.78 \pm 2.29$ & $24.48 \pm 2.27$ \\
        16 & 8 & $24.26 \pm 2.46$ & $22.85 \pm 2.25$ & $22.69 \pm 2.24$ & $23.93 \pm 2.17$ \\
        16 & 32 & $22.35 \pm 2.38$ & $24.23 \pm 2.24$ & $23.97 \pm 2.45$ & $22.77 \pm 2.37$ \\
        \addlinespace
        60 & Baseline & \multicolumn{4}{c}{$23.23\pm 2.02$} \\
        60 & 4 & $22.07 \pm 2.18$ & $23.73 \pm 2.40$ & $24.56 \pm 2.15$ & $20.45 \pm 2.16$ \\
        60 & 8 & $22.18 \pm 1.85$ & $24.85 \pm 2.28$ & $23.01 \pm 2.33$ & $23.38 \pm 1.50$ \\
        60 & 32 & $22.94 \pm 2.25$ & $18.68 \pm 1.45$ & $23.47 \pm 2.60$ & $23.10 \pm 2.03$ \\
        \bottomrule
    \end{tabular}
\end{table}


% @Andrew: result description + takeaway (which hypothesis does this extension validates?)
The temporal encoder extension is examined over three hyper parameter: history 
length ($K = 4, 8, 32$), GRU hidden dimension size ($H = 32, 64$), and actor-critic feature sharing 
(separated versus shared temporal extractors). Each hyperparameter combination is tested on both 
the fiducial 16-feature state and the extended 60-feature states for a total of 24 extension 
experiments.

Table~\ref{tab:temporal-eval} reports the held-out evaluation score for the temporal encoding 
extension hyperparameter sweep. The extension provides clearest benefit in the fiducial
16-feature observation configurations. The best model ($K = 4$, $H = 64$, separate) 
improves the evaluation score from $20.14\pm 2.25$ to $24.78\pm 2.29$. 
In contrast, the 60-feature configurations 
shows limited improvements from temporal encoding. Across hyperparameter variations, the 
temporal encoding extension yield less systematic effects than the choice of observation
feature set. 

In the fiducial 16-feature setting, the policy likely benefited from 
recent historical information encoded by GRU. The extended 60-feature setting, 
each state already contains historical summary features, making
the additional GRU summary partly redundant. In addition, the larger augmented 
representation may be harder for the actor-critic network of PPO to use effectively,
which could explain the underperforming 60-feature variants underperforms the baseline.

\subsection{Sample Reweighting}

The sample-reweighted PPO achieves a period-27 sweep score of \textbf{24.28}, compared to
20.14 for the uniform-sampling baseline which is a gain of 1.0 point (+4.0\%). The improvement is
consistent across training: the weighted training curve lies above the uniform baseline
throughout the full 10M environment steps, not only at convergence (Figure~\ref{fig:samplereweight-curve}).
This suggests that the gradient is being allocated more efficiently across the period
distribution, rather than the policy simply converging to a different local optimum.
 
% TODO: uncomment once figure is added
\begin{figure}[h]
  \centering
  \includegraphics[width=0.7\textwidth]{figures/sample_reweight_training.png}
  \caption{PPO training curves: uniform sampling (baseline) vs.\ score-weighted
  period sampling ($\alpha=0.9$, $T=30$). Weighted sampling outperforms
  throughout training.}
  \label{fig:samplereweight-curve}
\end{figure}
 
The result validates the hypothesis that \emph{where} the policy trains matters, not only
\emph{how} it trains. By concentrating rollouts on high-scoring pairs, the reweighted sampler lets the policy learn more from demonstrations of effective bidding behavior, i.e. correct pacing, CPA-compliant budget allocation, and selective bid
placement, rather than spending equal gradient budget on trajectories that exhibit poor strategies. [OPTIONAL, CAN BE REMOVED: The gain is modest compared to some other extensions
(e.g., market-regime conditioning at 29.7), which is expected: sampling changes the
\emph{efficiency} of learning but not the \emph{information} available to the policy.
The two interventions are complementary and could in principle be stacked.]

[CAN BE PUT INTO 4. DISCUSSION SECTION] One possible explanation for the modest 4\% gain is the aggressive choice of
$\alpha{=}0.9$, which places 90\% of sampling weight on the softmax distribution. This may cause the policy to overfit to high-scoring campaigns while under-training on harder advertisers with tighter CPA constraints or unusual budget profiles. A lower
$\alpha$ (e.g., 0.5--0.7) would retain more uniform coverage and may yield a better
balance between learning from good behavior and maintaining robustness across diverse
advertiser types. We leave this sweep to future work.



\subsection{Market-aware Feature Engineering}
@Bonnie: result description+ takeaway (which hypothesis does this extension validates?)

\subsection{Reward Shaping}
The reward shaping PPO achieves a best sweep score of $25.6$, at $\lambda = .0001$ and an improvement from the baseline. The gap for $\lambda = .01$ shows a visible gap in mean episode reward, which reflects the penalty term in the shaped objective and is not a drop in the unshaped score. We bounded sweep $\lambda$ at $.01$ as the penalty punishes overspend, and we do not have an additional reward mechanism for efficient bids that are within CPA and converged. A larger $\lambda$ would bias the policy toward conservative spending, which could hurt conversions that agents would otherwise win.
\begin{figure}[h]
  \centering
  \includegraphics[width=\textwidth]{figures/compare_obs16_default_vs_lambda_sweep.png}
  \caption{Reward shaping curves across $\lambda$ values versus baseline PPO. Training results are within similar convergence with the baseline.}
  \label{fig:reward_shaping}
\end{figure}
For each unique $\lambda$, the agent solved a different problem than original. We see $\lambda \in \{.0001, .001\}$ is too small to register in training, while $\lambda = .01$ drops training reward by $25\%$ without altering final reward in training. Exploratory runs of stricter $\lambda$ values, such as $.1 \text{ or } 1$, produced a policy toward ineffective bidding.
\subsection{RL training on LLM}

\begin{figure}[h]
    \centering
    \includegraphics[width=0.8\linewidth]{figures/llm_grpo_training.pdf}
    \caption{\textbf{GRPO + LoRA fine-tuning of Qwen3-8B on AuctionNet periods 7--26
    (validation on period 27).} Left: per-step training score (rolling-10 mean and
    min/max band over each batch of 16 rollouts), greedy validation reward at every
    checkpoint, and the frozen-base-model validation baseline (gray dotted line at
    17.08). Right: actor token entropy on a log scale, tracking how quickly the
    action distribution concentrates during training.}
    \label{fig:llm_grpo_training}
\end{figure}

We report the GRPO training of Qwen3-8B in
Figure~\ref{fig:llm_grpo_training}. The run uses rollout temperature \(1.0\),
group size \(G=4\), and no entropy bonus. During training, chain-of-thought generation was disabled due to the computational resource limit. The
validation curve shows the sample efficiency advantage of starting from a
pretrained policy: the frozen Qwen3-8B model already reaches a nontrivial
validation score of 17.08, and GRPO lifts it to 21 in the early stage
of training. The resulting score is below the peak of the best non-LLM
actor-critic policies, but the early jump suggests that pretraining gives the
policy a useful prior for converting auction states into bid multipliers.

The later training dynamics point to the limitation of this setup. With
chain-of-thought disabled, token entropy mainly reflects the spread of the
emitted action; once the policy concentrates around a narrow set of bid
multipliers, the \(G=4\) samples provide little within-group return variation
for GRPO. The current setup therefore reaches a useful low-entropy pacing mode
quickly, but larger rollout groups or stronger exploration pressure may be
needed to turn the sample-efficiency gain into competitive peak performance.


\section{Discussion and Limitations}
Several extensions showed better performance than original PPO. The principal limitation is computational resources, since each PPO run trains 10 million steps on the environment. Given the time constraint of the project, hyperparameter tuning is expensive. Future directions include finer hyperparameter tuning and mixed strategy of combining extensions to see if performance gain can be aggregated.
\paragraph{Reward Shaping}
{\footnotesize
% \paragraph{Acknowledgements.} We thank XXX.

\bibliographystyle{apalike}
\bibliography{references}
}

\section{Appendix (Optional)}

Include any details that don't fit within the main body here.

\end{document}
